\subsection{不变埃尔米特型}

在本节中，字母 k 表示域 R 或 C。令 V 为有限维 k-向量空间，\(\Phi\) 为 V 上分离的\(^{2}\)正埃尔米特型，G 为群，g 为 R-李代数，\(\rho: G \to \mathbf{GL}(V)\) 为群同态，\(\varphi: \mathfrak{g} \to \mathfrak{gl}(V)\) 为 R-李代数同态。

\begin{enumerate}
\def\labelenumi{\alph{enumi})}
\tightlist
\item
  形式 \(\Phi\) 在 G（分别地，g）下不变当且仅当 \(\rho(g)\) 关于 \(\Phi\) 是酉的（对所有 \(g \in G\)）（分别地，\(\varphi(x)\) 是反埃尔米特的\(^{3}\)（关于 \(\Phi\)，对所有 \(x \in g\)））。事实上，以 \(a^{*}\) 表示 V 的自同态 a 关于 \(\Phi\) 的伴随；对 G 中的 g、g 中的 x、V 中的 u 与 v，我们有
\end{enumerate}

\[
\begin{array}{c} \Phi (\rho (g) u, \rho (g) v) = \Phi (\rho (g) ^ {*} \rho (g) u, v), \\ \Phi (\varphi (x) u, v) + \Phi (u, \varphi (x) v) = \Phi ((\varphi (x) + \varphi (x) ^ {*}). u, v); \end{array}
\]

于是，\(\Phi(\rho(g)u,\rho(g)v)=\Phi(u,v)\)（对所有 \(u,v\)（在 V 中））当且仅当 \(\rho(g)^{*}\rho(g)=\mathrm{Id}_{\mathrm{V}}\)；类似地，\(\Phi(\varphi(x)u,v)+\Phi(u,\varphi(x)v)=0\)（对所有 \(u,v\)（在 V 中））当且仅当 \(\varphi(x)+\varphi(x)^{*}=0\)，由此即得所述断言。

\begin{enumerate}
\def\labelenumi{\alph{enumi})}
\setcounter{enumi}{1}
\item
  若形式 \(\varPhi\) 在 G（分别地，g）下不变，则 V 的稳定子空间的正交补是稳定的；特别地，表示 \(\rho\)（分别地，\(\varphi\)）此时是半单的（参见《代数学》第IX章）；此外，对所有 \(g \in G\)（分别地，\(x \in g\)），V 的自同态 \(\rho(g)\)（分别地，\(\varphi(x)\)）此时是半单的，其特征值的绝对值为 1（分别地，特征值都是纯虚数）；事实上，\(\rho(g)\) 是酉的（分别地，\(i\varphi(x)\) 是埃尔米特的，参见《代数学》第IX章）。
\item
  假设 \(k = \mathbf{R}\)。若 \(G\) 是连通李群，\(\rho\) 是李群态射，\(\mathfrak{g}\) 是 \(G\) 的李代数，\(\varphi\) 是由 \(\rho\) 诱导的同态，则 \(\varPhi\) 在 \(G\) 下不变当且仅当它在 \(\mathfrak{g}\) 下不变（第III章§6第5节推论3）。
\item
  V 上存在在 G 下不变的分离的正埃尔米特形式，当且仅当子群 \(\rho(\mathrm{G})\)（\(\mathbf{GL}(\mathrm{V})\) 的）是相对紧的（《积分学》第VII章§3第1节命题1）。
\end{enumerate}

